% Generated draft: requires mathematical and source review.
\documentclass[11pt]{article}
\usepackage[margin=25mm]{geometry}
\usepackage{fontspec}
\setmainfont{Latin Modern Roman}
\usepackage{amsmath,amssymb,mathtools,unicode-math}
\setmathfont{Latin Modern Math}
\usepackage{enumitem,needspace,xurl,hyperref}
\setlength{\parindent}{0pt}
\setlength{\parskip}{5pt}
\setlength{\emergencystretch}{3em}
\newcommand{\N}{\mathbb{N}}
\newcommand{\Z}{\mathbb{Z}}
\newcommand{\Q}{\mathbb{Q}}
\newcommand{\R}{\mathbb{R}}
\newcommand{\C}{\mathbb{C}}
\newcommand{\F}{\mathbb{F}}
\newcommand{\A}{\mathbb{A}}
\newcommand{\PP}{\mathbb{P}}
\newcommand{\E}{\mathbb{E}}
\newcommand{\eps}{\varepsilon}
\newcommand{\dd}{\,\mathrm{d}}
\begin{document}
\section*{UnsolvedMath ID: 30004582}
\textbf{Source problem code:} catalogue entry
\subsection{Definitions and mathematical statement}
\textbf{Terminology and scope.} The supplied record does not define ``everyday numeracy'' or specify the chronological boundaries of the ``long nineteenth century.'' It also does not specify a geographical setting, a population, or criteria for determining when someone has acquired everyday numeracy. These terms therefore remain unresolved historical qualifications; no numerical thresholds, dates, or population restrictions are imposed here.

\textbf{Source question.} The supplied statement is:
\begin{quote}
In the long nineteenth century, how did people acquire everyday numeracy?
\end{quote}

\textbf{Formulation limitation.} This is a historical research question about the acquisition of numeracy. The supplied record provides no mathematical proposition, quantified conjecture, or formal model whose truth is to be established. A precise mathematical statement cannot be extracted without adding assumptions and an intended mathematical target that the source does not supply. The question is retained with its unresolved scope for human review.
\subsection{Short English statement}
The question asks how people acquired everyday numeracy during the long nineteenth century. Its historical scope and its meaning of numeracy are unspecified. The supplied record labels it open, but does not formulate a mathematical conjecture.
\subsection{Sources}
\begin{itemize}
\item Local UnsolvedMath registry, \texttt{data/unsolvedmath/problems.json}, permanent ID \texttt{30004582}, source problem code \texttt{catalogue entry}; supplied title: \emph{Acquisition of Everyday Numeracy in the Nineteenth Century}. The record belongs to the Oberwolfach Reports Open Problems collection.
\item Supplied report citation: \emph{History of Mathematics: A Global Cultural Approach} (2021), \emph{Oberwolfach Reports}, pp.\ 2021--2092, DOI \texttt{10.4171/owr/2020/41}. Supplied source URL and literature reference labeled ``original Oberwolfach report'': \url{https://doi.org/10.4171/owr/2020/41}. The background describes Oberwolfach Reports as documenting mathematical workshops and their open-problem discussions.
\item Attribution metadata: the supplied proposed year is 2020; no individual proposer is named. The report citation gives publication year 2021.
\item Supplied literature assessment dated \texttt{2026-08-22}: status ``open.'' The assessment describes the acquisition of numeracy as a context-dependent historical research question and states that the extracted record contains no theorem or completed account justifying a stronger status. This is supplied dated curation, not an independently verified literature assessment. Its listed source is the original Oberwolfach report linked above; time-sensitive status and cited claims require independent verification.
\item The registry labels the statement ``exact'' and reports that it is internally readable and complete and was retained byte-for-byte. This supplied statement-verification note does not establish independent verification against the report.
\end{itemize}
\end{document}
